\documentclass{article}
\usepackage[slovene]{babel}

\usepackage{lipsum}

\usepackage{ragged2e}


\usepackage{lmodern}
\usepackage{graphicx}

\usepackage{enumitem}
\setlist[itemize,1]{noitemsep, topsep=0pt, label=\heartsuit}
\setlist[itemize,2]{noitemsep, topsep=0pt, label=\diamondsuit}
\setlist[enumerate,1]{label=\Alph*.}
\setlist[enumerate,2]{label*=\arabic*.}
\setlist[description]{style=nextline}

\usepackage{array,booktabs}

\usepackage{multirow}

\begin{document}


% \begin{flushleft}
%   Ta besedilo je 
%   levo poravnan.
%   \LaTeX{} se ne trudi narediti
%   vrstic z enakimi dolžinami.

%   \lipsum[1][1-2]
% \end{flushleft}

% \begin{FlushLeft}
%   Ta besedilo je 
%   levo poravnan.
%   \LaTeX{} se ne trudi narediti
%   vrstic z enakimi dolžinami.

%   \lipsum[1][1-2]
% \end{FlushLeft}

% \begin{flushright}
%   Ta besedilo je desno poravnano.
%   Kot prej tudi tu vrstice nimajo
%   enakih dolžin.

%   \lipsum[1][1-2]
% \end{flushright}

% \begin{FlushRight}
%   Ta besedilo je desno poravnano.
%   Kot prej tudi tu vrstice nimajo
%   enakih dolžin.

%   \lipsum[1][1-2]
% \end{FlushRight}

% \begin{Center}
%   V središču 
%   sveta.

%   \lipsum[1][1-2]
% \end{Center}

% \begin{center}
%   V središču 
%   sveta.

%   \lipsum[1][1-2]
% \end{center}

% \vspace{2cm}

% \section{Seznami}

% \begin{enumerate}
%   \item Različna okolja lahko mešamo
%   po lastnem okusu:
%     \begin{itemize}
%       \item Toda to lahko postane smešno.
%       \item[-] To se začne s pomišljajem.
%     \end{itemize}
%   \item Zapomnite si:
%     \begin{description}
%       \item[Neumne] zadeve ne bodo postale pametne,
%       če so v seznamu.
%       \item[Pametne] zadeve, za razliko, pa lahko čudovito
%       prikažemo s seznamom.
%     \end{description}
% \end{enumerate}

% \vspace{2cm}
% % enumitem

% Z uporabo možnosti \texttt{label=\textbackslash alph*)} lahko uporabljamo črke namesto številk:
% \begin{enumerate}%[label=\alph*)]
%   \item Z možnostjo \texttt{label=\textbackslash textit\{(\textbackslash roman*)\}}, dobimo \textit{(i)}, \textit{(ii)}, \textit{(iii)} \dots
%   \item Podobno, z možnostjo \texttt{label=\textbackslash textbf\{\textbackslash Alph*.\}}, dobimo \textbf{A.}, \textbf{B.}, \textbf{C.}...
% \end{enumerate}

% Lahko spremenimo oznake neurejenih seznamov.
% \begin{itemize}%[label=\heartsuit]
%   \item Tukaj uporabljamo \texttt{label=\textbackslash heartsuit}.
%   \item Z možnostjo \texttt{label=\textbackslash textendash} dobimo pomišljaj.
%   \item Z možnostjo \texttt{label=\textbackslash textasteriskcentered} dobimo zvezdico.
% \end{itemize}

% Z uporabo možnosti \texttt{style=nextline} lahko dosežemo, da se opisi začnejo v novi vrstici:
% \begin{description}%[style=nextline]
%   \item[Jabolko] Rdeče ali zeleno sadje.
%   \item[Pomaranča] Okroglo oranžno sadje.
% \end{description}

% %%%

% \begin{enumerate}%[label=\arabic*.]
%   \item  Prva točka na prvi stopnji.
%     \begin{enumerate}%[label*=\Alph*.]
%       \item Prva točka na drugi stopnji.
%       \item Druga točka na drugi stopnji.
%     \end{enumerate}
% \end{enumerate}

% \begin{enumerate}%[label=\Roman*., resume]
%   \item To je drugačen seznam, z rimskimi številkami.
%   \item Številke pa se nadaljujejo iz prejšnjega seznama.
%   \item Nikoli ne smete nadaljevati seznama z različnimi oznakami.
% \end{enumerate}

% \vspace{2cm}

% V preambuli:
% \usepackage{enumitem}
% \setlist[itemize,1]{noitemsep, topsep=0pt, label=\heartsuit}
% \setlist[itemize,2]{noitemsep, topsep=0pt, label=\diamondsuit}
% \setlist[enumerate,1]{label=\Alph*.}
% \setlist[enumerate,2]{label*=\arabic*.}
% \setlist[description]{style=nextline}


% V telesu dokumenta:
% \begin{itemize}
%   \item Neurejeni seznami prve stopnje imajo kot oznake srčke.
%   \item Opazujte, kako je bil razmik nastavljen tudi globalno.
%   \begin{itemize}
%     \item Neurejeni seznami na drugi stopnji imajo diamanti.
%   \end{itemize}
% \end{itemize}

% \begin{enumerate}
%   \item  Prav tako smo nastavili oznake urejenih seznamov.
%     \begin{enumerate}
%       \item Zdaj tisti na drugi stopnji privzeto podedujejo oznake (e.g. A.1.)
%     \end{enumerate}
% \end{enumerate}

% \begin{description}
%   \item[Opisni seznami] so nastavljeni tako, da se opisi začnejo v novi vrstici.
%   \item[To je] zelo uporabno, kadar so opisi daljši.
% \end{description}

% \section{Tabele}

Tabele v \LaTeX u se lahko poravnajo na sredino:
\begin{tabular}{ll} 1&2\\3&4 \end{tabular}
na dno:
\begin{tabular}[b]{ll} 1&2\\3&4 \end{tabular}
ali na vrh:
\begin{tabular}[t]{ll} 1&2\\3&4 \end{tabular}
glede na besedilo okoli njih.

\vspace{2cm}

\begin{tabular}{ l c r }
  \texttt{l} & \texttt{c} & \texttt{r} \\
  levo  &
  središče &
  desno  \\
  1 & 2 & 3 \\
\end{tabular}

\vspace{2cm}

\begin{tabular}{ l p{2cm} m{2cm} b{2cm}}
  \hline
  \texttt{l} & \texttt{p\{2cm\}} & \texttt{m\{2cm\}} & \texttt{b\{2cm\}} \\
  \hline
  Celica &
  Vrh besedila je poravnan s celico. &
  Središče besedila je poravnan s celico. &
  Dno besedila je poravnan s celico. \\\hline
\end{tabular}

\vspace{2cm}

\begin{tabular}{l|c r|}
  levo& središče & desno \\ \hline
  1 & 2 & 3 \\ \cline{1-1}
  4 & 5 & 6 \\ \cline{2-3}
  7 & 8 & 9 \\ \hline
\end{tabular}

\vspace{2cm}

\begin{tabular}
  { l @{\hspace{1cm}}  r @{,} l @{ €}}
  \toprule
  Izdelek & \multicolumn{2}{c}{Cena}\\
  \midrule
  Espresso & 1&50 \\
  Kava z mlekom & 2&50 \\
  Čaj & 1&80 \\
  \bottomrule
\end{tabular}

\end{document}